%% {{CODE}} chapter template -- {{PUBNAME}}
%%
%% For chapters in edited volumes and papers in proceedings. Remove `lineno`
%% from the options for the final version if the volume editor asks.
\documentclass[chapter,lineno,{{CITEOPT}}]{nextgeniq}

\publication{{{PUBNAME}}}{{{CODE}}}
\volumetitle{Title of the Volume}
\shorttitle{A short running title of up to 50 characters}

\title{Title of the Chapter: Concise, Specific and Informative}

\author[1]{First Author}
\author[2]{Second Author}
\affil[1]{Department, University, City, Country}
\affil[2]{Department, Institution, City, Country}
\correspondence{First Author, \texttt{first.author@example.org}}

\keywords{keyword one; keyword two; keyword three; keyword four}

\begin{document}

\maketitle

\begin{abstract}
Summarise the chapter's aim, approach and main contribution in 150--250 words.
Do not cite references in the abstract.
\end{abstract}

\section{Introduction}
\label{sec:intro}

Introduce the problem and explain how the chapter contributes to the theme of
the volume. Cite sources as \citet{smith2024} or \citep{lee2023,garcia2022}.
Explain specialised terms for readers from other disciplines.

\section{Background}
\label{sec:background}

Review the work the chapter builds on \citep{garcia2022}.

\section{Approach}
\label{sec:approach}

Describe and justify the methods. Number displayed equations so they can be
referred to, as in Equation~\eqref{eq:example}.
\begin{equation}
  \hat{y} = \arg\max_{y \in \mathcal{Y}} \; P(y \mid x; \theta)
  \label{eq:example}
\end{equation}

\subsection{A subsection}
Use subsections where they help the reader; avoid going deeper than three
levels.

\section{Findings}
\label{sec:findings}

Present the evidence, referring to every table and figure in the text, for
example Table~\ref{tab:results} and Figure~\ref{fig:example}.

\begin{table}[ht]
  \centering
  \caption{Table captions go above the table. Use horizontal rules only.}
  \label{tab:results}
  \begin{tabular}{lccc}
    \toprule
    Method        & Precision & Recall & F1 \\
    \midrule
    Baseline      & 0.81      & 0.77   & 0.79 \\
    Proposed      & 0.88      & 0.85   & 0.86 \\
    \bottomrule
  \end{tabular}
\end{table}

\begin{figure}[ht]
  \centering
  \includegraphics[width=0.6\linewidth]{example-image}
  \caption{Figure captions go below the figure. Supply figures at 300\,dpi or
  as vector graphics. Permission is needed for figures taken from other works.}
  \label{fig:example}
\end{figure}

\section{Discussion and Conclusion}
\label{sec:conclusion}

Discuss the implications and limitations, and suggest directions for further
research.

\begin{declarations}
  \declaration{Acknowledgements} Thank people and organisations who helped
  with the work.
  \declaration{Funding} Name the funders and grant numbers, or state that the
  research received no external funding.
  \declaration{Conflicts of interest} Declare any competing interests, or
  state that there are none.
  \declaration{Ethics approval} Give the approving body and reference where
  required, or state that approval was not required.
  \declaration{Use of generative AI} Describe any substantive use of
  generative AI tools, or state that none was used.
\end{declarations}

\bibliography{references}

\end{document}
